%! root=./main.tex
\section{Background and Motivation}~\label{sec:back}

We motivate the problem with a class of bugs that elude even experienced developers and resist exhaustive dynamic testing, then define our design scope and threat model.

\subsection{Motivation: ARM Unpredictable Behavior}

A CFI binary rewriter performs inherently complex transformations on executable code.
These include rewriting control-flow instructions (\eg branches), relocating and
reassembling both original and newly generated code segments, and embedding auxiliary
structures such as policy tables. As a result, producing a rewritten binary that is
both compliant with the underlying ISA specification \emph{and} faithful to CFI
safety requirements is nontrivial.

In particular, a CFI-enforcing binary rewriter must guarantee that every generated
instruction sequence remains fully compliant with the target ISA\@. However,
establishing such assurance through dynamic testing alone is fundamentally limited.

Like other ISAs, the ARM architecture explicitly defines \emph{unpredictable}
behaviors: execution outcomes that are left deliberately unspecified. These arise
from design choices intended to preserve implementation flexibility, to avoid
over-specifying invalid or ill-formed programs, and to balance performance with
backward compatibility. From a verification perspective, unpredictable behavior
introduces significant challenges: it induces semantic nondeterminism, necessitates
conservative over-approximation, and complicates the establishment of precise
correctness guarantees.

In practice, many CPU implementations handle such cases by trapping or terminating
execution, though the specification mandates no single outcome. Accordingly,
compilers and toolchains typically avoid generating instruction sequences that may
trigger unpredictable behavior. \pname{} explicitly disallows any rewritten binary
that could trigger such behavior.

This requirement motivates the use of formal verification in \pname{}: it ensures
that rewriting never introduces execution paths that enter architecturally
unpredictable regions. Achieving this guarantee is beyond the reach of conventional
dynamic testing, which explores only observed paths and cannot rule out problematic
sequences by construction. In contrast, formal verification enables proactive
identification of all such paths, allowing Cornerstone defense against control hijacking attacks, achieving widespread deployment with strong guarantees and near native overhead. However, the community has focused on design level guarantees, runtime overhead and policy generation largely overlooking implementation correctness and relying instead on dynamic testing. A more fundamental approach, formally verifying CFI implementations against safety and security specifications, remains largely unexplored. We need to be careful here to identify the overlook problem as developing guaranteed correct CFI implementations, not merely the problem of finding bugs in deployed CFI implementations. The letter is extensively studied in the literature, the former is not. This gap underscores a. them to be ruled out or corrected
statically. In Section~\ref{sec:case_studies}, we present two concrete error cases (branching
to the PC register and unaligned PC loads) that we encountered and corrected
through the formal verification process.


\subsection{Formally Verified Binary Rewriting}

\noindent\textbf{Positioning and scope.}
\pname{} is, to our knowledge, the first work to formally verify a \emph{program
transformation}, specifically a static binary rewriter, that takes arbitrary
binaries as input and produces security-hardened outputs that provably comply with
a precise safety specification.

\noindent\textbf{Problem setting.}
We study the formal verification of a system that operates on \emph{arbitrary
binaries}. In this work, we focus on arbitrary ELF binaries; within each binary,
we restrict our analysis to the \texttt{.text} section, \ie the executable
instruction stream. Other binary components, such as metadata, relocation entries,
and dynamic linking structures, may affect overall program behavior but are treated
as out of scope.

\noindent\textbf{Challenges of the domain.}
Formal methods for programs that consume and transform arbitrary binaries remain
rare and challenging. The input space is vast and highly irregular; reasoning
requires precise, low-level ISA semantics; and the system must preserve correctness
while actively modifying executable code, not merely analyzing or validating it.

\noindent\textbf{Relation to prior work.}
We contrast \pname{} with two representative lines of work.
\emph{CompCert}~\cite{compcert} is a seminal verified C compiler that
proves semantic preservation from source to target code using machine-checked
proofs. It replaces testing-based confidence with formal guarantees and produces
efficient code. However, it assumes well-defined source programs and operates on
source-level inputs rather than arbitrary binaries. \pname{} works directly on
existing binaries and verifies the correctness of a \emph{rewriting transformation},
rather than compilation from a high-level language.
\emph{RockSalt}~\cite{rocksalt2012pldi} verifies a validator that enforces
software fault isolation (SFI) for x86 binaries, guaranteeing that any accepted
binary satisfies the sandboxing policy. However, RockSalt does not verify the
\emph{rewriter} that produces such binaries; it assumes untrusted code and checks
it at load time. \pname{} complements this approach by verifying the transformation
itself, ensuring that generated binaries satisfy security properties by construction.

\noindent\textbf{Choice of target ISA.}
We target the ARMv7 ISA for three reasons. First, ARMv7 represents one of the most
widely deployed architectures by installed base, particularly in low-powered embedded
and IoT devices that are commonly deployed without modern security defenses, making
binary hardening for this platform practically important. Second, a formal
specification for the ARMv7 instruction decoder exists and has been the subject of
substantial prior work on formal verification of instruction decoding, providing a
solid semantic foundation. Third, by targeting a fixed-width instruction encoding,
we avoid the variable-length decoding complexities of x86, allowing us to focus
formal reasoning effort on the binary rewriter itself rather than the decoder.


\subsection{Threat Model and Research Scope}\label{sec:threat_model}
\kjee{NOTE: a typical threat model text for CFI paper. This is still an *optional*.}
\ixb{Tricky things to include ($\sim$21:30): basically attackers get to send inputs to victim programs and send malware to the CFI rewriter, out of scope is messing with the hardware, ELF metadata, and manipulating the policy input.}

Unlike prior work that verifies either compilers (\ie source-to-binary
transformations) or validators (\ie binary-to-accept/reject decisions),
\pname{} verifies a \emph{binary-to-binary transformation} that enforces security
properties directly on arbitrary inputs.

\pname targets adversaries who exploit memory corruption vulnerabilities to hijack program control flow, a threat model consistent with prior work on code-reuse attacks. We assume the attacker can corrupt code pointers (return addresses, function pointers, or vtable entries) and possesses a writable memory region under their control (\eg a heap object or auxiliary buffer). We further assume the attacker can perform limited memory disclosure sufficient to infer code and data addresses in the presence of ASLR\@.
Under \pname's CFI enforcement, any control-flow transfer that violates the policy is redirected to \texttt{abort()}, preventing attacker-controlled diversion of execution.

The current implementation targets statically linked, non-PIE ARM binaries, in which all control transfers are self-contained and symbol
references use absolute addresses. Extending coverage to PIE and dynamically linked binaries is achievable with additional engineering effort, and we leave it for future work.
